\section{Apple: asset balances do not identify asset movements}
\label{sec:apple-operating-assets}

The preceding methods need receipts, lots, useful lives and settlement dates.
Apple's disclosures supply some of the balances those events must explain,
but they do not supply the events themselves. Its FIFO and straight-line
policies narrow the permitted cost calculations without identifying the actual
inventory lots or asset cohorts \citep[Note 1, printed p.~34]{apple2024}.
Comparing the notes with the cash-flow statement shows why filling those
missing inputs by subtracting balances can produce a wrong history.

\subsection{An inventory adjustment is not purchases or payments}

Between 28 September and 28 December 2024, components fall from 3,627 to 2,792
million dollars while finished goods rise from 3,659 to 4,119 million. Total
inventory falls from 7,286 to 6,911: the component changes, $-835$ and $+460$,
sum to $-375$ million \citep[Note 5, printed p.~9]{apple2025q1}. Those two
classes do not reveal quantities, unit costs or the transfers and sales that
produced the balances.

The quarter's cash-flow statement reports $+215$ million for inventories within
its reconciliation of net income to operating cash \citep[printed p.~5]{apple2025q1}.
That is an indirect reconciliation adjustment, not total cash paid to suppliers.
Even replacing it with minus the stock change would give $+375$, a different
quantity. The difference $215-375=-160$ million remains unexplained by these
selected disclosures. In fiscal 2024, inventory rose from 6,331 to 7,286, while
the cash-flow adjustment was $-1{,}046$; relative to the $-955$ stock-change
shortcut the difference is $-91$ million \citep[printed pp.~31 and 33]{apple2024}.
Neither residual is an identified currency effect, write-down or other journal
entry. Production, consumption, ownership, valuation and payment histories must
be established before they can explain the differences.

\subsection{Testing a proposed fixed-asset history}

The fixed-asset notes distinguish three measurements. The table shows gross
property, plant and equipment (PP\&E), its accumulated depreciation and the
remaining carrying amount, all in millions of dollars
\citep[Note 5, printed p.~39]{apple2024}
\citep[Note 5, printed p.~9]{apple2025q1}.

\begin{center}
\begin{tabularx}{\linewidth}{@{}Xrrr@{}}
\toprule
Amount & 30 Sep 2023 & 28 Sep 2024 & 28 Dec 2024\\\midrule
Gross PP\&E & 114,599 & 119,128 & 120,615\\
Accumulated depreciation & 70,884 & 73,448 & 74,546\\
Net PP\&E & 43,715 & 45,680 & 46,069\\\bottomrule
\end{tabularx}
\end{center}

Accumulated depreciation is shown here as a positive deduction: gross less
that deduction equals net. During fiscal 2024, gross cost increases 4,529,
accumulated depreciation increases 2,564 and net PP\&E increases 1,965.
Yet the note reports PP\&E depreciation expense of 8.2 billion, and the
cash-flow statement reports capital payments of 9,447 million. Its broader
depreciation-and-amortisation adjustment is 11,445 million
\citep[printed pp.~33 and 39]{apple2024}. These amounts differ in meaning;
the 8.2 billion also has coarser displayed precision. The change in accumulated
depreciation cannot be substituted for the year's depreciation charge.

Consider a deliberately limited asset book to test a tempting reconstruction.
Let $G_0,G_1$ be opening and closing gross cost and $B_0,B_1$ accumulated
depreciation, recorded as positive amounts. Let $A$ be gross additions, $R$
the original cost of retired assets, $d$ the depreciation charge and $S$ the
accumulated depreciation removed with those assets. Assume there are no other
movements and all depreciation is expensed. For this illustration, use the
displayed amounts as exact, including $d=8{,}200$ million. The two balances
must then satisfy
\begin{equation}
 G_1=G_0+A-R,\qquad B_1=B_0+d-S,\qquad 0\leq S\leq R.
 \label{eq:apple-asset-movement}
\end{equation}
The inequality requires a nonnegative carrying amount for the retired assets
in this cost model. If additions are equated with the 9,447 of cash payments,
the gross balance requires $R=9{,}447-4{,}529=4{,}918$. The depreciation
balance requires $S=8{,}200-2{,}564=5{,}636$. The proposed book removes more
accumulated depreciation than original cost, by 718 million. It is infeasible
under these combined assumptions. That finding rejects the proposed history;
it does not establish an error in Apple's filing or identify which assumption
must change.

The result is not sensitive to small rounding differences under a specified
rounding model. If the million-dollar balances and payment are rounded to the
nearest million, and 8.2 billion to the nearest 0.1 billion, the excess
$d-9{,}447+(45{,}680-43{,}715)$ lies in the closed outer envelope
$[666.5,769.5]$ million. Its lower endpoint is
$8{,}150-9{,}447.5+45{,}679.5-43{,}715.5=666.5$.
This is a conditional precision check, not a verified issuer rounding rule;
the exact tie convention remains unknown.

Allowing additions to be partly unpaid provides two simple repairs. Retain
the same selected opening and closing balances, 8,200 depreciation and 9,447
cash payments. Assume retirement at the start of the year with zero proceeds,
depreciation of surviving assets, and new acquisitions at year end. There are
no opening capital payables, taxes or other asset movements. The figures below
are constructed histories, not estimates of actual Apple transactions.

\begin{center}
\begin{tabularx}{\linewidth}{@{}Xrr@{}}
\toprule
Annual movement (USD million) & History 1 & History 2\\\midrule
Gross additions & 10,529 & 11,529\\
Gross cost retired & 6,000 & 7,000\\
Accumulated depreciation removed & 5,636 & 5,636\\
Retired carrying amount, expensed & 364 & 1,364\\
Cash paid for additions & 9,447 & 9,447\\
Additions remaining payable & 1,082 & 2,082\\\bottomrule
\end{tabularx}
\end{center}

Both histories satisfy equation~\eqref{eq:apple-asset-movement}; their gross
additions can also be allocated to match the three reported asset classes.
For History 1, acquisition debits PP\&E cost 10,529 and credits cash 9,447 and
a capital payable 1,082. Depreciation debits expense and credits accumulated
depreciation by 8,200. Retirement debits accumulated depreciation 5,636 and
loss 364, while crediting gross PP\&E cost 6,000. Net PP\&E rises 1,965 and
cash falls 9,447, so total assets fall 7,482. Liabilities rise 1,082 and equity
falls 8,564 through the two expenses: $1{,}082-8{,}564=-7{,}482$.
History 2 has the same asset change but liabilities rise 2,082 and equity
falls 9,564. These are coherent subsystem consequences, not a fit to all of
Apple's statements or an exhaustive family of real asset histories. Neither
history receives a probability or preference from this arithmetic.

The following quarter supplies a new snapshot, not the missing annual movements.
Its gross, accumulated and net changes are 1,487, 1,098 and 389 million,
whereas capital payments are 2,940 and broad depreciation/amortisation is
3,080 \citep[printed pp.~5 and 9]{apple2025q1}. The inspected interim note does
not separately identify PP\&E depreciation expense. The lease note adds a
related distinction: finance-lease right-of-use assets are included in PP\&E,
while operating-lease right-of-use assets are in other noncurrent assets.
The 1.0 billion of noncash right-of-use acquisitions covers the disclosed lease
population \citep[Note 8, printed p.~42]{apple2024}; adding it all to PP\&E
or to cash capital expenditure would erase that distinction.

The method therefore preserves reported balances, cash payments, expense and
unresolved movement components separately. A future cost or funding query
still needs lot and asset ages, useful lives, residual values, impairment and
settlement facts. The real filing has supplied constraints and ruled out one
simple history, while leaving several different explanations open.

\subsection{Local alternatives need different company histories}

The two asset books expose a constraint on combining explanations. Both fit
the selected PP\&E movements and capital payments. Can they therefore be
inserted interchangeably into one otherwise unchanged history of Apple?
The journals above answer no. Relative to History 1, History 2 retires an
additional 1,000 of gross cost without removing additional accumulated
depreciation, so it recognises an additional 1,000 loss. It acquires an
additional 1,000 of gross cost on credit. Those changes cancel in net PP\&E
and leave cash unchanged, but they add 1,000 to liabilities and subtract
1,000 from profit and equity. Amounts remain millions of dollars, under the
same conditional asset-book assumptions, including exact illustrative
depreciation of 8,200 and no tax effect within that subsystem.

Let $P_i$ denote the vector of statement effects generated by asset History
$i$, and let $R_i$ denote the effects of the rest of the company's history.
For this comparison, order each vector as asset change, liability change,
profit, equity change and cash change. Both full histories must explain the
same reported vector $Y$. If an admissible remainder exists for each branch,
their necessary relation is
\[
P_1+R_1=Y=P_2+R_2,
\qquad
R_2-R_1=-(P_2-P_1)=(0,-1{,}000,+1{,}000,+1{,}000,0).
\]
Thus the rest of History 2 must contribute 1,000 less liability growth and
1,000 more profit and equity growth, with no additional asset or cash change.
Holding the remainder fixed is wrong relative to the common-statement target.
The balance-sheet identity alone fails to detect the error: adding 1,000 to
liabilities and subtracting 1,000 from equity still leaves their sum unchanged.

The reported annual totals make the required differences concrete. Apple's
asset growth is 12,397, liability growth is 17,593, profit is 93,736 and equity
change is $-5{,}196$ \citep[printed pp.~29--32]{apple2024}. Since History 1
contributes liability growth of 1,082, profit of $-8{,}564$ and equity change
of $-8{,}564$, its hypothetical remainder must contribute 16,511, 102,300
and 3,368 respectively. History 2 instead requires 15,511, 103,300 and
4,368. In both cases the remaining asset change must be 19,879. These are
necessary subtraction results. They do not supply the transactions, expense
classifications, taxes or liability assignments that would make a remainder
admissible.

The distinction matters for cash as well. In the isolated asset book,
depreciation and the retirement loss are noncash expenses. Adding both back
to its negative profit gives zero operating cash; acquisitions use 9,447 of
investing cash. Switching branches therefore adds 1,000 to the noncash loss
and 1,000 to the corresponding addback, leaving subsystem cash unchanged.
If that loss were mapped into the annual cash-flow statement's broad Other
adjustment of $-2{,}266$, the remaining contributions to that row would have
to be $-2{,}630$ in History 1 and $-3{,}630$ in History 2
\citep[printed p.~33]{apple2024}. That mapping is itself conditional: the
source does not identify either retirement loss or the composition needed
to offset it. A matching operating-cash total cannot substitute for the
reported adjustment rows.

Nor may a remainder erase the cash-population difference from
Section~\ref{sec:apple-statements}. After the common asset payments, the
required remaining balance-sheet cash change is 9,425, whereas the required
remaining cash-flow change is 8,653. Their 772 difference is inherited from
the disclosed opening populations. Calling it another cash receipt would
manufacture a transaction rather than explain the record.

A proposed completion must now identify the liability class containing the
capital payable, the expense class containing the loss, the tax consequences
and the noncash cash-flow mapping, and satisfy the detailed statement rows
under the same assignments. It must also retain the asset cohorts and payment
terms needed at the next date. The December snapshot does not establish that
either hypothetical annual payable remained outstanding or identify when it
was settled. Failure of the unchanged-remainder construction therefore calls
for a coupled historical model. It rejects that construction without deciding
whether either local asset book has an admissible full-company completion.


\subsection{An offset must explain the individual rows}
\label{sec:apple-coupled-completion}

The required offset can be tested without pretending that its transaction has
been observed. Suppose a separate economic obligation would justify an
expense and an unpaid liability of $q$ in History 1 and $q-1{,}000$ in
History 2, with $q\geq1{,}000$. These are alternative amounts in two histories;
they are not a release booked within either history. For a numerical example,
take $q=2{,}000$: each branch debits expense and credits its own liability,
for 2,000 or 1,000 respectively. The difference supplies the 1,000 lower
liability and 1,000 higher profit required of History 2's remainder.
The economic obligation is an additional premise. Subtracting reported
totals has not established that it exists.

Even conditional on that premise, a matched total is insufficient. Assume
both the retirement loss and the new expense enter product cost, and both
the capital payable and the new liability enter the same non-tax current
liability category. History 2's extra loss then offsets its smaller expense
within that expense row, and its larger capital payable offsets its smaller
expense liability within that liability row. Other operations, distributions,
share counts and OCI remain fixed in this comparison. Equal tax expense and
tax balances are a further condition: equal pretax profit does not imply
equal deductions, tax bases or deferred tax.

The indirect cash-flow statement exposes another restriction. If the
expense liability is treated as an ordinary operating accrual, its 1,000
smaller increase reduces the working-liability adjustment by 1,000. The
extra retirement loss increases the noncash addback by 1,000. Total operating
cash flow is unchanged, but the separately reported Other adjustment and
Other current and non-current liabilities adjustment each change. That
construction fails to reproduce Apple's detailed cash-flow statement,
even though it balances and matches total operating cash
\citep[printed p.~33]{apple2024}.

A different conditional reporting map identifies the expense liability's
recognition as a separate noncash adjustment. Its amount must then be
excluded from the working-liability adjustment to avoid adding it back
twice. Let $\alpha$ be the fraction of this hypothetical recognition
allocated to the separate noncash adjustment; the remaining $1-\alpha$
enters the operating-liability adjustment. This parameter describes an
assumed allocation, not a probability or a freely selectable accounting
policy. With the retirement loss also assigned to Other, the differences
between History 2 and History 1 are
\[
\Delta\mathrm{Other}=1{,}000-\alpha\,1{,}000,
\qquad
\Delta\mathrm{WorkingLiabilities}=-(1-\alpha)\,1{,}000,
\qquad
\Delta\mathrm{CFO}=0.
\]
For $\alpha=0$, the two rows differ by $+1{,}000$ and $-1{,}000$.
For $\alpha=\tfrac12$, they still differ by $+500$ and $-500$.
Both row differences vanish only at $\alpha=1$ within this construction.
Including the full liability movement as well as the separate addback
would instead understate the History 2 minus History 1 operating-cash
difference by 1,000.

Under the common expense and liability mappings, equal taxes and the
disjoint noncash mapping with $\alpha=1$, all five annual statements are
unchanged by the combined transformation. Fixed shares also preserve EPS;
fixed OCI and distributions carry the profit equality into comprehensive
income and the equity statement. This is a conditional invariance result.
It assumes a feasible full-company History 1 exists and that the transformed
events remain admissible. It neither constructs that base nor proves that
either complete history is compatible with all notes. Assigning the two
expenses to different rows, the liabilities to different categories, or
different tax consequences immediately breaks the relevant equality.

The notes make these remaining conditions specific. Note 5 separates gross
PP\&E by asset class: the extra acquisition and retirement must cancel
within each affected class to leave that table unchanged. Note 6 separates
26,601 of current income taxes payable from 51,703 of other current
liabilities, and 9,254 of non-current income taxes payable from 36,634 of
other non-current liabilities. A total-liability offset cannot choose its
own location inside those amounts. Note 8 identifies finance-lease
right-of-use assets of 1,069 within PP\&E and distinguishes lease liabilities
and noncash lease acquisitions. Those disclosures do not identify the
hypothetical capital purchases or expense liability. Nor do Note 12's
purchase commitments and contingency assessment establish the assumed
obligation or its amount
\citep[printed pp.~39, 42, 45--46]{apple2024}.

The distinction therefore survives the repair: an algebraic offset is
available; a row-preserving transformation is available under stated
economic, tax and reporting assumptions; an actual Apple completion is
unestablished. To admit one, the historical model needs the underlying
obligation, gross asset classes, expense and liability assignments, tax
bases and noncash treatment, together with a feasible rest of the company.
The subsequent quarter additionally needs settlement dates and surviving
asset cohorts. Identical annual statements do not make those future
responses identical.


\subsection{Equal liabilities can require different cash paths}
\label{sec:apple-settlement-paths}

The conditional repair leaves different obligations beneath the same annual
liability total. In the $q=2{,}000$ illustration, History 1 owes 1,082 for
capital purchases and 2,000 for the separate expense; History 2 owes 2,082
and 1,000 respectively. Each owes 3,082. Consider what paying those balances
would do, conditional on the economic and reporting assumptions needed for
the preceding construction. Assume payment at carrying value, sufficient
unrestricted cash, and no additional interest, foreign exchange,
remeasurement or tax effect. Capital settlement is classified as investing
cash and expense settlement as operating cash in this example. Neither
the obligation nor its settlement terms have been identified in Apple's
disclosures.

If both obligations are paid at the first subsequent date, each payment
debits its liability and credits cash. No new expense arises: the expense
or acquisition was recognised in the prior period. Both histories reduce
cash and liabilities by 3,082, leave assets other than cash unchanged, and
leave profit, OCI and equity unchanged. The indirect operating-cash
calculation includes the expense liability's decrease once; it excludes
capital settlement from the operating-liability adjustment. The
cash-flow statement nevertheless distinguishes the histories. Writing
each ordered pair as operating cash followed by investing cash, the
joint alternatives are
\[
(\mathrm{CFO},\mathrm{CFI})_{H_1}=(-2{,}000,-1{,}082),
\qquad
(\mathrm{CFO},\mathrm{CFI})_{H_2}=(-1{,}000,-2{,}082).
\]
Both pairs sum to $-3{,}082$. Selecting the smaller outflow from each
category would instead produce $(-1{,}000,-1{,}082)$, whose total is only
$-2{,}082$. No branch produces that pair. Selecting both larger outflows
would invent an extra 1,000 payment. The category intervals can describe
marginal possibilities, but their Cartesian product contains paths
excluded by the shared history.

Timing makes the distinction relevant even to total cash. Suppose capital
obligations are paid at the first date and expense obligations at the
second. History 1 then pays 1,082 followed by 2,000; History 2 pays 2,082
followed by 1,000. Relative to History 1, History 2 has 1,000 less cash
after the first date and the same cumulative cash after the second.
Paying the expense obligations first reverses the first-date difference.
An aggregate liability of 3,082 suffices for the narrow query of total
undiscounted cash after both payments, under these fixed assumptions.
It does not suffice for cash by date or by statement category. Retaining
liability kind, amount, settlement terms and joint branch identity is
therefore necessary for those queries.

The actual following quarter constrains a proposed continuation without
identifying it. Apple reports PP\&E payments of 2,940, operating cash of
29,935, investing cash of 9,792 and total cash growth of 356 in the quarter
ended 28 December 2024 \citep[printed p.~5]{apple2025q1}. If both
hypothetical obligations settled in that quarter, the remaining PP\&E
payments would have to be 1,858 under History 1 or 858 under History 2.
The rest of the company would have to contribute operating cash of
31,935 or 30,935, and investing cash of 10,874 or 11,874, respectively.
After unchanged financing cash of $-39{,}371$, either required remainder
has cash growth of 3,438; subtracting the hypothetical 3,082 settlement
recovers the reported 356. These are necessary conditions on unconstructed
remainders. Positive remaining payments and a matching cash identity do
not establish the underlying transactions or prove that either
continuation fits the complete quarterly record.

This calculation extends the conditional history into explicit journals
and statement effects while keeping actual-history admission unresolved.
If cash is insufficient, the model must retain the resulting funding or
default question; it cannot invent a loan to preserve the path. A later
depreciation or production query requires different information again:
asset cohorts, remaining lives, residual values, capacity and tax bases.
The settlement example supplies none of those quantities. A change in the
query must therefore reopen the state-sufficiency check, even when the
annual statements and the final cash total coincide.

